\clearpage
\onecolumn
\section{Hypotheses and Failure-Mode Interpretation}
\label{sec:supp-theory}

This section records theory and interpretation. The evidence is suggestive
rather than mechanism-complete, so the discussion stays at hypothesis level.

\begin{center}
\centering
\includegraphics[width=\textwidth,height=0.74\textheight,keepaspectratio]{figures/fig_supplement_theory_evidence_bridge.png}
\captionsetup{hypcap=false}
\captionof{figure}{Theory bridge linking the hypothesis-level interpretation to
render evidence. The bands connect material salience,
target-view grounding contracts, embodied-route sensitivity, and a
selected-failure render wall.}
\label{fig:supp-theory-evidence-bridge}
\end{center}
\clearpage

\begin{center}
\centering
\includegraphics[width=0.98\textwidth,height=0.74\textheight,keepaspectratio]{figures/fig_supplement_theory_failure_mode_map.png}
\captionsetup{hypcap=false}
\captionof{figure}{Failure-mode interpretation map for the hypothesis-level
section. Each lane anchors one interpretation to tracked render evidence and
states the scope of the example.}
\label{fig:supp-theory-failure-mode-map}
\end{center}
\clearpage

\begin{center}
\centering
\includegraphics[width=0.86\textwidth,height=0.72\textheight,keepaspectratio]{figures/fig_supplement_theory_hypothesis_boundary_companion.png}
\captionsetup{hypcap=false}
\captionof{figure}{Hypothesis-scope companion. Render rows keep
material, grounding, navigation, selected failure render wall, and
cross-hypothesis evidence visible. The render-backed lens summarizes how the
examples should be read.}
\label{fig:supp-theory-hypothesis-scope-companion}
\end{center}
\clearpage

\twocolumn

\subsection{Material Salience Hypothesis}

The material salience hypothesis is that conversion changes matter most when a
target is recognized through localized material cues rather than silhouette
alone. A cup, bottle, or backpack can remain geometrically visible while losing
opacity, emission, normal, displacement, clearcoat, or procedural texture cues.
If a VLM relies on those cues to separate the target from nearby objects, a
small material perturbation can produce a grounding delta even when global
image similarity remains high.

\subsection{Coordinate-Format Hypothesis}

The coordinate-format hypothesis is that material conversion and coordinate
protocol failures can interact. A material shift can change visual salience,
while a model-specific coordinate convention can change the returned point
frame. If the analysis reports only raw point hits, a target-centered camera can
hide this interaction. Separating answer accuracy, raw point hits,
normalized-1000 point hits, parser coverage, and pair agreement is therefore
part of the causal interpretation rather than bookkeeping.

\subsection{Embodied-Stack Sensitivity Hypothesis}

The navigation stack sees more than a static target crop. It receives
egocentric observations, map context, route memory, stopping decisions, and
failure thresholds. A converted scene can be static-render comparable while a
selected rollout changes local route shape, endpoint behavior, or final
navigation error. Conversely, a visible material difference can be irrelevant
to the selected route. The 99-episode paired run is included to test whether
converted scenes can be exercised by a real official stack, not to prove broad
embodied robustness.

\subsection{Why Selected Failures Still Matter}

Selected failures provide mechanism visibility rather than rates. A selected
clearcoat target-loss case warns that static conversion
success can hide target preservation failures. A selected procedural-texture
case warns that checker or generated texture evidence may disappear during
conversion. A selected navigation divergence warns that route inspection is
useful even when aggregate metrics are limited to the selected run. These examples help define
future tests: they say what to test before estimating how often it fails.

\subsection{Negative-Result Discipline}

The supplement keeps negative cases visible. Procedural texture remains a
known limitation case. NVIDIA clearcoat target loss remains a selected baseline
failure, and Qwen coordinate behavior remains a coordinate-protocol effect.
This discipline keeps the ACL contribution focused on language-grounded scene
conversion measurement rather than an asset-converter leaderboard.
